\chapter{Round Robin, EF1 and EFX}\label{ch:background}

FUNCTION: This chapter provides the concepts that every later result relies on and explains why
round robin achieves EF1 but not necessarily EFX.



\section{Setting and Notation}\label{sec:setting}
FUNCTION: Introduces agents, goods, additive valuations, allocations, picking sequences and
round robin, together with the number of full rounds and the remainder of goods. All later
chapters use this notation.

The Round Robin algorithm and the related fair division problem are considering a set of agents $N = \{1,...,n\}$ as well as a set of undividable goods $M = \{g_1,...,g_m\}$. Each agent has an additive \textit{valuation function} $v_i(S)$, whereas $S \subseteq M$ and $v_i(S):\sum_{g\in S} v_i(g)$, with $v_i(g) \geq 0$. The model in \textit{see ch. 3} generates strict positive, continuos values. 
An \textit{allocation} $A = \{A_1,...,A_n \}$ is a tuple of subsets of M, so that every agent $i \in N$ gets a \textit{bundle}, $A_i \subseteq M$, $A_i \cap A_j = \emptyset$ for every par of agents $i, j \in N$ and $\bigcup_{i \in N} A_i = M$, assigned. That means $A$ is has disjunct subsets and is complete (every good gets allocated) $A = M$.

I define a \textit{picking sequence} as a sequence of the length $m$ each number represents one Agent $i$. After \cite{bouveret2011}: each agents picks his most valuable available good in his turn.

The central picking sequence of this work, the Round Robin algortihm, has the sequence §1,2,...,n,1,2,...,. divided by m turns. I define one round as Round $t$, which contains the turns $(t-1)n+1$ until $tn$. The last round can be incomplete. The good that Agent $i$ takes in Round $t$ is $i_t$. So the set goods for Agent i is the allocation $A_i$ for each Agent $i$.

For Further usage I follow \cite{manurangsi2021} and write to variables $r$ and $q$. r is the amount of \textit{full rounds} defined as $r = \lfloor \frac{m}{n} \rfloor$ and, $q = m - nr$ which is the rest of goods in the last round. Applied to the Round Robin algorithm, that means that the Agents $1,...,q$ get $r+1$ goods and the other remaining agents $r$ goods.

\subsubsection{Example}
Let $n=3$, $m=7$, thus $r = \lfloor \frac{7}{3} =2 \rfloor$ and $q = 1$: the sequence of Round Robin is 1231231. The bundle sizes for the agents are 3, 2, 2



\section{Fairness Notions}\label{sec:notions}
FUNCTION: Defines envy freeness, EF1, EFX and its approximate version, the alpha-EFX ratio used
as the target quantity, each attributed to its source. The section shows how the notions relate
to each other.

\section{Why Round Robin Is EF1, and What EFX Additionally Requires}\label{sec:buffer}
FUNCTION: Restates the known argument that round robin is EF1 and reads it as a buffer that is
never negative. EFX holds exactly when this buffer covers a gap between two single goods. A
worst case example in which the buffer vanishes while the gap stays large shows that round robin
has no constant approximation guarantee for EFX. This motivates the average case question of the
thesis. The buffer and the gap are the explanatory tool for all results in Chapter 4.

\section{Does an EFX Allocation Exist?}\label{sec:existence}
FUNCTION: Summarises where the literature proves EFX existence and defines the region of our
design where it is open. For that region the chapter states explicitly that we assume existence,
and gives the reasons for this assumption. This tells the reader for every result whether a
failure of round robin is a failure of the mechanism or possibly of the problem itself.

